\documentclass[10pt]{article}

%% Language and font encodings
\usepackage[english]{babel}
\usepackage[utf8x]{inputenc}
\usepackage[T1]{fontenc}
\usepackage{url}

%% Sets page size and margins
\usepackage[a4paper, top=2cm, bottom=2cm, left=2cm, right=2cm, marginparwidth=1.75cm]{geometry}

%% Useful packages
\usepackage{graphicx}

\begin{document}
\begin{center}

    \includegraphics[scale=0.5]{assets/logo.png}

    \begin{LARGE}
        \textbf{SCAR: A Blockchain based approach for academic registry}
    \end{LARGE}

    \vspace{.75cm} Projeto e Seminário

    \vspace{.25cm} Licenciatura em Engenharia Informática e Computadores

    \vspace{.75cm}
    \begin{large}
        \begin{tabular}{ c c c }
            Diogo Rodrigues               & Gonçalo Frutuoso              \\
            \texttt{49513@alunos.isel.pt} & \texttt{49495@alunos.isel.pt} \\
        \end{tabular}

        \vspace{.75cm} Orientadores \vspace{0.25cm}

        \begin{tabular}{ c c }
            Cátia Vaz                    & Alexandre Francisco      \\
            \texttt{cvaz@cc.isel.ipl.pt} & \texttt{aplf@tecnico.pt}
        \end{tabular}
    \end{large}

    \vspace{.75cm} \today \vspace{.75cm}

\end{center}

\section{Introdução}

A veracidade e autenticidade dos registos académicos, especialmente certificados e certidões de notas,
têm sido um desafio persistente no campo da educação. A confiabilidade desses documentos é crucial não apenas
para os próprios indivíduos, mas também para instituições de ensino, empregadores e outras entidades que
dependem desses registros para tomar decisões importantes. No entanto, a atual gestão desses registros
ainda é amplamente baseada em métodos tradicionais, como impressões em papel ou documentos digitais simples,
o que os torna vulneráveis à adulteração, falsificação ou perda. A crescente incidência de fraudes
relacionadas a certificados acadêmicos, juntamente com os desafios associados à validação manual desses
documentos, destaca a necessidade urgente de uma solução inovadora e segura.

Atualmente, embora existam algumas soluções digitais disponíveis para a gestão de registros acadêmicos,
muitas delas carecem de uma abordagem verdadeiramente segura e confiável. Os sistemas de gestão académicos
(AMS \cite{LinWays}) e os portais de estudantes geralmente oferecem opções para visualizar e baixar certificados e certidões
de notas, porém, esses documentos digitais ainda podem ser facilmente manipulados ou falsificados. Além disso,
a verificação desses registros é muitas vezes um processo manual e demorado, envolvendo a comunicação direta entre
as partes interessadas ou a consulta a bases de dados centralizadas, o que pode comprometer a eficiência
e a integridade do processo.

Diante desse cenário, surge a necessidade premente de uma solução que possa garantir a autenticidade e a
imutabilidade dos registros acadêmicos, ao mesmo tempo em que simplifica e agiliza o processo de validação.
É nesse contexto que nossa proposta de projeto, denominada SCAR (Smart Contracts for Academic Registry),
se destaca como uma abordagem inovadora e eficaz para a gestão de certificados e certidões de notas por
meio da tecnologia blockchain e smart contracts.

FALTA MAIS UM PARAGRAFO A INDICAR OS OBJETIVOS DO TRABALHO

\section*{Tecnologias}

EXPLICAR PORQUE VAMOS FALAR DESTAS Tecnologias
Embora, não tenha sido lecionado ao longo da licenciatura considerou-se interessante o desafio de aprender
as mesmas pois irão dar uma solução multiplataforma.(A MELHORAR)

\subsubsection*{React Native}

React Native \cite{ReactNativeBook}, juntamente com a plataforma Expo \cite{Expo}, emerge como uma das ferramentas mais versáteis
para o desenvolvimento de aplicações móveis. Ao unir a eficiência do React.js com a capacidade de compilação
nativa, React Native permite aos programadores criar aplicações móveis robustos e de alto desempenho
com uma base de código compartilhada. A plataforma Expo, por sua vez, oferece um conjunto de ferramentas
e serviços adicionais que simplificam e aceleram o processo de desenvolvimento, permitindo aos desenvolvedores
concentrarem-se na criação de experiências de utilizador excepcionais.

\subsubsection*{Solidity}

Solidity, a linguagem de programação utilizada para escrever smart contracts na blockchain Ethereum,
o valor do tempo investido pelo nosso grupo é ainda mais evidente. À medida que a adoção da tecnologia
blockchain continua a crescer e se expandir para além do mundo das criptomoedas, o domínio de Solidity se
torna uma habilidade altamente valorizada e procurada. Smart Contracts desempenham um papel fundamental em
uma ampla gama de aplicações distribuidas (dApps), desde finanças descentralizadas (DeFi) até votação eletrônica
e muito mais.

\section{Requisitos}

ADICIONAR AQUI UM TEXTINHO

\subsection*{Intervenientes}

A solução SCAR envolve três principais intervenientes: os estudantes ou alumni, os empregadores ou entidades terceiras e os administradores
de faculdade. Cada um desempenha um papel crucial no ecossistema da plataforma.

\subsection*{Papel de cada interveniente e ações que podem realizar}

\begin{itemize}

    \item \textbf{Estudantes ou Alumni:} Os candidatos são a essência da plataforma SCAR. Eles terão a capacidade de registrar-se na
          aplicação móvel, onde poderão carregar e armazenar seus certificados acadêmicos de forma segura na blockchain.
          Além disso, os candidatos poderão visualizar e compartilhar seus certificados com empregadores durante processos
          de recrutamento, proporcionando uma experiência transparente e eficiente.

    \item \textbf{Administradores de Faculdade:} Os administradores da faculdade são responsáveis por fornecer e validar
          os certificados acadêmicos emitidos pela instituição de ensino. Terão a capacidade de acessar a plataforma SCAR
          para emitir e autenticar os certificados dos candidatos, garantindo assim a integridade e a validade dos documentos.

    \item \textbf{Empregadores ou Entidades externas:} Os empregadores desempenham um papel fundamental ao solicitar a verificação de certificados
          acadêmicos durante entrevistas de emprego. Ao acessar a plataforma SCAR, os empregadores podem verificar
          instantaneamente a autenticidade dos certificados apresentados pelos candidatos, garantindo assim um processo
          de recrutamento mais confiável e transparente.

\end{itemize}

\subsection*{Requisitos obrigatórios}

\begin{itemize}

    \item\textbf{Implementação da aplicação móvel utilizando React Native com a plataforma Expo:} A escolha do React Native
          juntamente com a plataforma Expo proporciona uma abordagem eficaz para o desenvolvimento de aplicações móveis
          multiplataforma. Isso garante uma experiência consistente para os utilizadores, independentemente do dispositivo
          utilizado, e facilita o processo de desenvolvimento para a equipa de desenvolvimento.

    \item\textbf{Utilização de smart contracts em Solidity para armazenar e validar os certificados acadêmicos na blockchain
              Ethereum:} A utilização de smart contracts em Solidity oferece uma solução segura e confiável para armazenar
          e validar os certificados acadêmicos na blockchain Ethereum. Essa abordagem garante a integridade e a
          imutabilidade dos dados, tornando a plataforma SCAR altamente confiável e resistente a fraudes.

    \item\textbf{Desenvolvimento de funcionalidades para registo, autenticação e armazenamento seguro de certificados na
              blockchain:} O desenvolvimento de funcionalidades robustas para registro, autenticação e armazenamento seguro
          de certificados na blockchain é essencial para garantir a segurança e a confiabilidade da plataforma SCAR.
          Essas funcionalidades devem ser projetadas com foco na usabilidade e na segurança, proporcionando uma
          experiência intuitiva e transparente para os utilizadores.

\end{itemize}

\subsection*{Requisitos opcionais}

\begin{itemize}

    \item\textbf{Integração de funcionalidades adicionais, como notificações em tempo real e compartilhamento de certificados
              através de meios digitais:} A integração de funcionalidades adicionais, como notificações em tempo real e
          compartilhamento de certificados através de meios digitais como por exemplo redes sociais, pode melhorar ainda mais a experiência do
          utilizador na plataforma SCAR, aumentando a usabilidade e adesão de novos utilizadores.

\end{itemize}

\pagebreak

\section{Calendarização}

\begin{table}[!h]
    \centering
    \begin{tabular}{| c | c | c |}
        \hline
        Data de início  & \begin{tabular}{@{}c@{}}Duração \\ (semanas)\end{tabular} & Descrição                                                        \\
        \hline
        16 de Fevereiro & 2                                                         & Estudo sobre a framework de frontEnd a utilizar,                 \\ & & \textbf{Expo React Native}                  \\
        \hline
        28 de Fevereiro & 2                                                         & Estudo sobre a tecnologia Smart Contracts                        \\ & & e a linguagem que esta suporta, \textbf{Solidity}        \\
        \hline
        6 de Março      & 2                                                         & Elaboração da proposta.                                          \\
        \hline
        18 de Março     & 2                                                         & Decisão acerca da estrutura e modelo relacional a utilizar.      \\
        \hline
        29 de Março     & 1                                                         & Elaboração do mockup do frontend da aplicação                    \\
        \hline
        5 de Abril      & 2                                                         & Implementação dos ecrãs da aplicação                             \\
        \hline
        5 de Abril      & 3                                                         & Implementação dos \textbf{Smart Contracts} necessários.          \\
        \hline
        15 de Abril     & 1                                                         & Elaboração do relatório de progresso e apresentação individual.  \\
        \hline
        30 de Abril     & 2                                                         & Testes, avaliação experimental incluindo testes de usabilidade.  \\
        \hline
        13 de Maio      & 2                                                         & Criação e desenvolvimento do cartaz e entrega da versão beta.    \\
        \hline
        28 de Maio      & 6                                                         & Otimizações, finalização do relatório e entrega da versão final. \\
        \hline
    \end{tabular}
    \caption{Calendarização do projeto.}
    \label{calendar}
\end{table}



Ao longo do trabalho efetuado semanalmente, irá ser elaborado continuadamente documentação, testes unitários e relatório do projeto, de forma a que no final do mesmo, o relatório esteja completo e pronto a ser entregue.

\bibliographystyle{abbrv}
\bibliography{references}

\end{document}